\documentclass[10pt, conference, letterpaper]{IEEEtran}

\usepackage{algorithm}
\usepackage{algorithmicx}
\usepackage{algpseudocode}
\usepackage{amsfonts}
\usepackage{amsmath}
\usepackage{amssymb}
\usepackage[ansinew]{inputenc} 
\usepackage{xcolor}
\usepackage{mathtools}
\usepackage{graphicx}
\usepackage{caption}
\usepackage{subcaption}
\usepackage{import}
\usepackage{multirow}
\usepackage{cite}
\usepackage[export]{adjustbox}
\usepackage{breqn}
\usepackage{mathrsfs}
\usepackage{acronym}
%\usepackage[keeplastbox]{flushend}
\usepackage{setspace}
\usepackage{bm}
\usepackage{stackengine}

\usepackage{listings}

\lstset{%
 backgroundcolor=\color[gray]{.85},
 basicstyle=\small\ttfamily,
 breaklines = true,
 keywordstyle=\color{red!75},
 columns=fullflexible,
}%

\lstdefinelanguage{BibTeX}
  {keywords={%
      @article,@book,@collectedbook,@conference,@electronic,@ieeetranbstctl,%
      @inbook,@incollectedbook,@incollection,@injournal,@inproceedings,%
      @manual,@mastersthesis,@misc,@patent,@periodical,@phdthesis,@preamble,%
      @proceedings,@standard,@string,@techreport,@unpublished%
      },
   comment=[l][\itshape]{@comment},
   sensitive=false,
  }

\usepackage{listings}

% listings settings from classicthesis package by
% Andr\'{e} Miede
\lstset{language=[LaTeX]Tex,%C++,
    keywordstyle=\color{RoyalBlue},%\bfseries,
    basicstyle=\small\ttfamily,
    %identifierstyle=\color{NavyBlue},
    commentstyle=\color{Green}\ttfamily,
    stringstyle=\rmfamily,
    numbers=none,%left,%
    numberstyle=\scriptsize,%\tiny
    stepnumber=5,
    numbersep=8pt,
    showstringspaces=false,
    breaklines=true,
    frameround=ftff,
    frame=single
    %frame=L
}

\renewcommand{\thetable}{\arabic{table}}
\renewcommand{\thesubtable}{\alph{subtable}}

\DeclareMathOperator*{\argmin}{arg\,min}
\DeclareMathOperator*{\argmax}{arg\,max}

\def\delequal{\mathrel{\ensurestackMath{\stackon[1pt]{=}{\scriptscriptstyle\Delta}}}}

\graphicspath{{./figures/}}
\setlength{\belowcaptionskip}{0mm}
\setlength{\textfloatsep}{8pt}

\newcommand{\eq}[1]{Eq.~\eqref{#1}}
\newcommand{\fig}[1]{Fig.~\ref{#1}}
\newcommand{\tab}[1]{Tab.~\ref{#1}}
\newcommand{\secref}[1]{Section~\ref{#1}}

\newcommand{\rev}[1]{{\color{blue}#1}}
\newcommand{\red}[1]{{\color{red}#1}}
\newcommand{\mytexttilde}{{\raise.17ex\hbox{$\scriptstyle\mathtt{\sim}$}}}

%\renewcommand{\baselinestretch}{0.98}
% \renewcommand{\bottomfraction}{0.8}
% \setlength{\abovecaptionskip}{0pt}
\setlength{\columnsep}{0.2in}

% \IEEEoverridecommandlockouts\IEEEpubid{\makebox[\columnwidth]{PUT COPYRIGHT NOTICE HERE \hfill} \hspace{\columnsep}\makebox[\columnwidth]{ }} 

\title{NNDL project report template}

\author{Jacopo Pegoraro$^\dag$
\thanks{$^\dag$Department of Information Engineering, University of Padova, \newline email: \texttt{\{name.surname\}@unipd.it}}
\thanks{Special thanks / acknowledgement go here.}
} 

\IEEEoverridecommandlockouts

\newcounter{remark}[section]
\newenvironment{remark}[1][]{\refstepcounter{remark}\par\medskip
   \textbf{Remark~\thesection.\theremark. #1} \rmfamily}{\medskip}

\begin{document}

\maketitle

\begin{abstract}
Social distancing and temperature screening have been widely employed to counteract the COVID-19 pandemic, sparking great interest from academia, industry and public administrations worldwide. 
While most solutions have dealt with these aspects separately, their combination would greatly benefit the continuous monitoring of public spaces and help trigger effective countermeasures.
This work presents milliTRACE-IR, a joint mmWave radar and infrared imaging sensing system performing unobtrusive and privacy preserving human body temperature screening and contact tracing in indoor spaces. milliTRACE-IR combines, via a robust sensor fusion approach, mmWave radars and infrared thermal cameras. 
It achieves fully automated measurement of distancing and body temperature, by jointly tracking the subjects's faces in the thermal camera image plane and the human motion in the radar reference system. 
Moreover, milliTRACE-IR performs contact tracing: a person with high body temperature is reliably detected by the thermal camera sensor and subsequently traced across a large indoor area in a non-invasive way by the radars. When entering a new room, a subject is re-identified among several other individuals by computing gait-related features from the radar reflections through a deep neural network and using a weighted extreme learning machine as the final re-identification tool.
Experimental results, obtained from a real implementation of milliTRACE-IR, demonstrate decimeter-level accuracy in distance/trajectory estimation, inter-personal distance estimation (effective for subjects getting as close as $\boldsymbol{0.2}$~m), and accurate temperature monitoring (max. errors of $\boldsymbol{0.5}~^{\circ}$C). Furthermore, milliTRACE-IR provides contact tracing through highly accurate ($\boldsymbol{95\%}$) person re-identification, in less than $\boldsymbol{20}$ seconds.\\ 

\rev{This is an example abstract from a paper we wrote a couple of years ago. I would say an abstract should not be longer than $250$ words. Here, you should briefly state: 
\begin{enumerate}
\item the technical scenario/field of research and its timeliness/relevance in general (one sentence),
\item what you do in the report/paper and why it is important, how it advances the state of the art in its field (a few sentences), 
\item summarize the main and best results of your study/proposal/method (one or two sentences),
\item (optional) how others could benefit from your results for further research, or within commercial products (one sentence).
\end{enumerate}  
The abstract is one of the most important parts of the paper/report. You have roughly one minute to catch the reader's attention. A poor abstract may already move you towards the rejection side in the reviewer's decision process. In the abstract, 1) establish the context, 2) motivate the problem, 3) briefly describe the solution, and 4) present the main results of your work. Ideally, use one (short) sentence for each of the previously mentioned items to keep your abstract short. Overall, this should be a short summary of the whole content of your paper, including your results.}\\

\rev{See the abstract as a personal challenge for each of your papers. Finally, the abstract should contain the main message about your work, so that the reader will know what she/he can find even without reading it (as is the case most of the time). The abstract is a mini-paper on its own and, as such, it is a major endeavor to write.}\\ 

\red{I suggest writing the Abstract as the very last thing. You may sketch it at the beginning, but then always finalize it at the end.}
\end{abstract}

\IEEEkeywords
Unsupervised Learning, Optimization, Neural Networks, Convolutional Neural Networks, Variational Autoencoders. \rev{A list of keywords defining the tools and the scenario. I would not go beyond {\it six} keywords.}
\endIEEEkeywords


\section{Introduction}
\label{sec:introduction}

Wi-Fi sensing infers human motion from variations of channel state information
(CSI), reusing commodity radios without recording identifiable images or audio.
Its practical value, however, depends on representations that separate motion
from environment-, device-, and packet-dependent distortions. SHARP addresses
this requirement by transforming CSI into phase-stable Doppler maps before
classification, and reports person- and environment-independent human activity
recognition (HAR) with commodity access points \cite{meneghello2023sharp}.
The transformation is not a plain Fourier transform: it estimates sparse paths,
selects a phase reference, sanitizes nuisance phase, and only then applies the
Doppler analysis.

This project asks whether the expensive part of that pipeline can be
approximated directly. Given a local window of raw complex CSI, a neural network
is trained to reproduce the Doppler map generated by SHARP. If successful, this
would preserve an established classifier interface while replacing iterative
signal processing with a fixed feed-forward computation. It also poses a useful
scientific test: can generic convolutional networks learn the phase invariance
that model-based preprocessing imposes explicitly?

The primary research question is:
\emph{Can an end-to-end neural network replace SHARP's optimization-intensive
CSI preprocessing while preserving both sparse motion structure and downstream
HAR performance?} Three supporting questions separate possible causes of
failure. First, do frequency-aware decoding, antenna-independent weights, exact
temporal alignment, and motion-aware training prevent a center-spectrum
shortcut? Second, do pixel reconstruction metrics predict whether generated
maps remain useful to a frozen SHARP classifier? Third, what does a comparison
with explicit affine phase correction reveal about the transformation that a
student must learn?

The contribution is a feasibility study rather than a new state-of-the-art HAR
method:
\begin{itemize}
    \item a reproducible CSI-to-SHARP-map distillation pipeline with controlled
    chronological splits and task-aware evaluation;
    \item a progressively corrected architecture and training setup addressing
    antenna sharing, subcarrier structure, temporal support, sparse motion, and
    Doppler locality;
    \item evidence that a full-resolution student partially preserves held-out
    motion and HAR utility, together with a class-level analysis of the
    remaining smoothing failure and a stronger explicit affine baseline.
\end{itemize}

\figref{fig:pipeline} summarizes the comparison. Both pixel fidelity and
classifier utility use the official precomputed SHARP representation: students
regress its maps, and each generated representation is sent through a classifier
trained on the same map distribution. The rest of the report reviews related
representations, defines the data and learning problem, presents the final model
and loss, and answers the research questions experimentally.

\begin{figure*}[t]
    \centering
    \includegraphics[width=\textwidth]{pipeline.pdf}
    \caption{Evaluation design. Raw CSI is transformed by the SHARP teacher, a
    direct neural student, or a simple affine phase correction followed by the
    same fixed Hann/STFT tail. Reconstruction metrics compare maps with the
    official SHARP target; a frozen classifier measures retained HAR utility.}
    \label{fig:pipeline}
\end{figure*}


\section{Related Work}
\label{sec:related}

\subsection{Invariant Wi-Fi Representations}

SHARP estimates a sparse channel impulse response and uses its strongest path as
a phase reference before constructing Doppler spectrograms
\cite{meneghello2023sharp}. Its representation is designed to transfer across
people and environments. Widar3.0 pursues the same general objective through a
physics-derived body-coordinate velocity profile, showing the value of removing
domain-dependent propagation effects before recognition
\cite{zheng2019widar}. Both systems therefore place substantial inductive bias
before the classifier rather than asking a generic classifier to discover
invariance from raw CSI.

\subsection{Phase-Aware CSI Processing}

Measured CSI contains amplitude variation and packet-dependent phase errors
caused by timing, carrier-frequency, sampling-frequency, and hardware offsets.
Such nuisance terms can dominate the much smaller phase changes induced by
motion. Ratnam \emph{et al.} formulate CSI preprocessing explicitly and analyze
how preprocessing choices affect sensing \cite{ratnam2024optimal}. The SHARP
dataset further documents the acquisition setting and its 80-MHz, 256-tone CSI,
of which 242 subcarriers are retained \cite{meneghello2023dataset}. These works
motivate treating phase correction as part of the learning problem, not as a
minor input normalization.

\subsection{Learning-Assisted Signal Processing}

Learning can reduce signal-processing cost without discarding known structure.
SLNet learns a spectrogram through a neural architecture tailored to wireless
sensing instead of regressing an unconstrained image
\cite{yang2023slnet}. LISTA unrolls sparse coding into a finite learned network,
preserving the optimization structure while learning efficient updates
\cite{gregor2010lista}. U-Net supplies the multiscale skip-connected backbone
used in our students \cite{ronneberger2015unet}, but it does not itself encode
complex phase equivariance or wireless propagation. The gap addressed here is
therefore narrower than generic end-to-end HAR: we test whether a convolutional
student can directly distill the output of a phase-aware preprocessing pipeline,
and compare that attempt with retaining a minimal explicit phase correction.


\section{Problem Formulation and Data Pipeline}
\label{sec:problem}

\subsection{Supervised Reconstruction Task}

For each local window, the input is raw CSI
\(\mathbf{X}\in\mathbb{C}^{A\times S\times T_x}\), with \(A=4\) receiver
streams, \(S=242\) usable subcarriers, and \(T_x=370\) packets. Real and
imaginary components are exposed as separate numeric channels, so a batch has
shape
\[
    [B,4,242,370,2].
\]
The target \(\mathbf{Y}\in[\floorval,1]^{B\times4\times340\times100}\)
contains one normalized 100-bin Doppler spectrum per antenna and output time.
A 31-packet analysis window requires 15 packets of context on either side, so
370 packets align exactly with 340 valid target frames. The objective is to
learn \(f_\theta(\mathbf{X})\approx\mathbf{Y}\) without exposing activity
labels to the student.

\subsection{Teacher and Reference Distributions}

The official precomputed targets follow the SHARP chain
\cite{meneghello2023sharp}: packet-level normalization, sparse path estimation,
selection of the strongest path as phase reference, affine phase sanitization,
Hann-windowed Doppler analysis, per-frame normalization, and flooring at
\(\floorval\approx0.0631\). We pair these paper-sourced maps with their raw CSI,
and use the same map distribution both for regression metrics and for the frozen
HAR classifier trained by the provided reproduction notebook.

The final experiment uses 92 temporally aligned recordings from all 12
available legacy AR scenario folders in the public 80-MHz CSI dataset
\cite{meneghello2023dataset}. Five of 97 candidate pairs are excluded because
their raw and Doppler lengths do not have SHARP's expected temporal offset;
forcing them into training would misalign input and target frames. Within every
remaining recording, the first \(80\%\) is available for training,
\(80\)--\(90\%\) for validation, and \(90\)--\(100\%\) for final testing. All
12 scenarios contribute to training and source validation; classifier-compatible
S1a/S1b/S1c determine early stopping and final reporting. A 31-frame guard
between intervals prevents context leakage. Training windows use stride 170,
while frozen-classifier evaluation uses stride 30 for denser coverage.
Uncompressed NumPy memory maps and deterministic recording-aware windowing make
this protocol reproducible without changing its samples.

\section{Learning Framework}
\label{sec:framework}

\subsection{Architecture Progression}

The initial temporal U-Net flattened antenna, subcarrier, and complex
components into channels, then assigned an independent output channel to every
antenna--Doppler pair. This makes a recording-average vertical spectrum an easy
solution. The next model shared one single-antenna network across all antennas
and added a spatial output head, but retained a 25-bin spectral bottleneck and
interpolation. The final model removes both shortcuts. \tabref{tab:models}
records the representative configurations and W\&B runs used in the report.

\begin{table*}[t]
\centering
\caption{Representative student architectures. Parameter counts are measured
from the saved configurations.}
\label{tab:models}
\footnotesize
\setlength{\tabcolsep}{4pt}
\begin{tabular}{@{}p{0.20\textwidth}p{0.47\textwidth}p{0.09\textwidth}p{0.16\textwidth}@{}}
\hline
Model (representative run) & Main inductive bias and limitation & Parameters & Training domain \\
\hline
Temporal 1-D U-Net (baseline) &
Frequency and antenna axes flattened into channels; independent temporal
output channels make the mean spectrum easy to represent. &
3.245M & PI, 30 carriers \\
Shared-antenna spatial head (\texttt{hpm4mjl4}) &
One antenna-independent temporal core and a 2-D head; retains a coarse
25-bin grid followed by spectral interpolation. &
0.793M & legacy AR, 242 carriers \\
Full-resolution shared 2-D (\texttt{az9k6ori}) &
Joint frequency--time encoder, GroupNorm, direct 100-bin grid, full 2-D
decoder, and exact valid-time crop. &
1.314M & legacy AR, 242 carriers \\
\hline
\end{tabular}
\end{table*}

\subsection{Full-Resolution Shared-Antenna Model}

The final architecture processes antennas independently with identical weights.
It reshapes
\[
[B,A,S,T,2]\rightarrow[BA,2,S,T],
\]
so antenna identity cannot become a shortcut, but every antenna--window pair is
retained. A \(9\times7\) frequency--time stem and two strided 2-D stages produce
feature tensors of sizes
\[
\begin{split}
[BA,64,61,370]&\rightarrow[BA,96,31,185]\\
&\rightarrow[BA,128,16,92].
\end{split}
\]
Residual blocks use GroupNorm, avoiding recording-composition-dependent
BatchNorm statistics. Only after these convolutions have modeled local
subcarrier structure is the remaining frequency axis averaged; a standard
\(1\times1\) projection then mixes latent channels.

The bottleneck is projected directly to
\([BA,16,92,100]\). A proper 2-D time--Doppler decoder upsamples time through
92, 185, and 370 positions and merges two projected encoder skips. It predicts
all 100 Doppler bins directly, with no 25-to-100 spectral interpolation. The
output head produces \([BA,370,100]\); removing exactly 15 frames from each
side yields the valid \([B,A,340,100]\) output. \figref{fig:architecture}
summarizes this flow. The model has 1,314,385 trainable parameters. These
changes correct the identified decoder and alignment flaws, but they do not
impose complex phase equivariance.

\begin{figure*}[t]
    \centering
    \includegraphics[width=\textwidth]{architecture.pdf}
    \caption{Tensor flow of the final student. Antennas share one network;
    subcarrier structure is processed before average pooling; the decoder
    predicts the complete time--Doppler grid and applies the exact 31-packet
    valid crop.}
    \label{fig:architecture}
\end{figure*}

\subsection{Motion-Aware Objective and Sampling}

Most target energy lies near zero Doppler, so plain MSE rewards a stationary
center ridge. We use
\begin{equation}
\mathcal{L}=\mathcal{L}_{\mathrm{map}}
 +0.25\mathcal{L}_{\mathrm{motion}}
 +0.25\mathcal{L}_{\mathrm{leak}}
 +0.05\mathcal{L}_{W}.
\label{eq:loss}
\end{equation}
\(\mathcal{L}_{\mathrm{map}}\) is MSE over all antennas, times, and 100 bins.
Bins 45--55 are designated stationary. Outside them,
\(\mathcal{L}_{\mathrm{motion}}\) weights squared error continuously by target
activity \(\operatorname{clip}((Y-\floorval)/(1-\floorval),0,1)\).
\(\mathcal{L}_{\mathrm{leak}}\) penalizes predicted power above the floor where
the target is at the floor. Finally, \(\mathcal{L}_{W}\) normalizes positive
off-center power into a distribution and averages the absolute difference
between predicted and target cumulative sums along the 100-bin Doppler axis.
This one-dimensional Wasserstein term distinguishes a nearby displaced peak
from equally wrong power far away.

Training uses deterministic motion stratification. Within each recording, the
top quartile of windows by active off-center frame fraction is marked
motion-rich. Every epoch includes all 1,748 rich windows and an equal-size,
non-duplicated rotating subset of ordinary windows. Batches contain 16 rich and
16 ordinary windows drawn across eight recordings. Validation and test remain
unmodified. The optimizer is Adam with learning rate \(10^{-3}\), automatic
mixed precision, a 50-epoch limit, and patience 8 on S1a/S1b/S1c validation
loss. Training completed all 50 epochs and restored the best checkpoint from
epoch 45. The smaller comparison model uses all stride-30 windows and batch
size 256.


\section{Experimental Evaluation and Discussion}
\label{sec:results}

\subsection{Protocol and Metrics}

The student is trained on \(0\)--\(80\%\) of each recording and validated on
\(80\)--\(90\%\). The frozen SHARP classifier was trained on S1a/S1b/S1c over
\(0\)--\(60\%\) and selected on \(60\)--\(80\%\); no classifier weights are
updated. Both are evaluated on the same disjoint \(90\)--\(100\%\) interval,
with stride 30 and 718 windows from the five classifier activities. This is
generalization to unseen temporal portions of known recordings, not to unseen
recordings, people, or environments. Each method receives identical raw
windows. The recording-mean centering expected by the classifier is estimated
independently for every representation and does not use labels.

We report full-map and off-center MSE against windowed targets from the official
precomputed SHARP Doppler traces. A frame is active when its maximum outside
bins 45--55 exceeds 0.2;
precision, recall, and F1 compare predicted and target support. The motion-mass
ratio is total predicted divided by target off-center excess power, with one
ideal. Frozen classifier accuracy is the primary task metric.

\subsection{RQ1: Can the Student Learn the Mapping?}

\textbf{Memorization sanity check.}
Before full-data training, two students were optimized repeatedly on the same
16 motion-rich PI windows. The 30-subcarrier model reached
\(3.881\times10^{-5}\) MSE, while the 242-subcarrier model reached
\(4.194\times10^{-5}\)(\figref{fig:diagnostics}(a)). Both reproduced the selected peaks. Thus, the end-to-end training path is functioning
correctly, and the model has sufficient capacity to memorize the selected local
examples, although this test says nothing about generalization.

\textbf{Full-data behavior.}
With official precomputed SHARP Doppler supervision, the final model behaves
very differently from the preliminary students. Motion-aware training loss decreases from
0.03397 to 0.01264 over all 50 epochs. All-scenario validation decreases from
0.02800 to 0.01415, while S1a/S1b/S1c validation decreases from 0.02506 to
0.01326 and reaches its minimum, 0.01212, at epoch 45 (\figref{fig:diagnostics}(b)). Both validation curves
continue improving, so there is no observed deterioration on the held-out
temporal intervals. Because training uses motion-balanced sampling while
validation retains the natural distribution, their absolute loss gap is not
interpreted.

\begin{figure*}[t]
    \centering
    \includegraphics[width=\textwidth]{diagnostics.pdf}
    \caption{Optimization diagnostics. (a) Both 30- and 242-subcarrier models
    memorize 16 fixed windows. (b) With official precomputed SHARP Doppler
    supervision, full-data
    train and validation objectives improve throughout training; the dashed
    line marks the selected epoch-45 checkpoint.}
    \label{fig:diagnostics}
\end{figure*}

The held-out example in \figref{fig:qualitative} shows that the model has learned an input-dependent mapping rather than producing a fixed average Doppler pattern. The student places off-center responses at the target event times. Its main error is now spatial smoothing:
thin target peaks become wider lobes, and nearby temporal events merge. The
complete corrected configuration avoids the center-only shortcut in this
example, but does not preserve all fine motion geometry.

\begin{figure*}[t]
    \centering
    \includegraphics[width=\textwidth]{qualitative.pdf}
    \caption{Held-out S1a\_C window 16920--17260 with identical color limits.
    The full-resolution student localizes the principal target events but
    broadens them; affine sanitization followed by the fixed STFT retains
    sharper time--Doppler structure.}
    \label{fig:qualitative}
\end{figure*}

\subsection{RQ2: Do Reconstructed Doppler Traces Preserve Downstream HAR
Performance Under a Frozen Classifier?}

\tabref{tab:main-results} gives the matched comparison. The antenna-shared
frequency--time U-Net reaches \(63.09\%\), a 19.78-point improvement over the
antenna-shared temporal U-Net with a coarse spectral head and well above raw
STFT. The complete corrected architecture and training
configuration therefore improves both map and downstream metrics.

\begin{table*}[t]
\centering
\caption{Held-out S1a/S1b/S1c comparison over the same 718 windows. All map
metrics use windowed targets from the official precomputed SHARP Doppler
traces.}
\label{tab:main-results}
\footnotesize
\setlength{\tabcolsep}{5pt}
\begin{tabular}{@{}lrrrrr@{}}
\hline
Representation & Accuracy (\%) & Balanced (\%) & Off-center MSE & Active F1 &
Mass ratio \\
\hline
Official precomputed SHARP Doppler & \textbf{91.36} & 91.74 & 0 & 1.000 & 1.00 \\
Raw fixed STFT & 21.73 & 20.00 & 0.128148 & 0.357 & 33.28 \\
Direct affine sanitization + fixed STFT & \textbf{91.36} & \textbf{91.80} & \textbf{0.000155} & \textbf{0.943} & \textbf{0.98} \\
Full-resolution frequency--time student & 63.09 & 58.86 & 0.001062 & 0.758 & 1.21 \\
Shared temporal student + coarse spectral head & 43.31 & 40.05 & 0.002758 & 0.505 & 2.22 \\
\hline
\end{tabular}
\end{table*}

The \(63.09\%\) aggregate accuracy does not imply classifier-equivalent or
uniform recognition: it is helped by \(100\%\) empty-room, \(99.36\%\) still,
and \(91.67\%\) running accuracy. Walking and jumping reach only \(1.30\%\) and
\(1.98\%\), respectively; 152 of 154 walking windows and 92 of 101 jumping
windows are classified as running. Balanced accuracy is consequently only
\(58.86\%\). The student has learned whether and approximately where motion
occurs, as supported by F1 \(0.758\) and mass ratio 1.21, but smoothing removes
the spectral and temporal details needed to distinguish dynamic activities.

In this provenance-matched experiment, off-center MSE and support F1 broadly
rank the methods in the same order as classifier accuracy. They remain
insufficient by themselves: active support treats a broad running-like lobe as
successful even when it erases class-discriminative shape. Reconstruction and
task metrics must therefore be reported together.

\subsection{RQ3: What Does a Simplified Phase-Sanitization Baseline Reveal?}

We include a non-learned baseline to test whether explicitly encoding a known
phase-error structure is sufficient, rather than asking the student to discover
it. The procedure is adapted from the affine phase-sanitization stage in
SHARP's public pipeline \cite{meneghello2023sharp}, consistent with established
CSI phase-error models \cite{ratnam2024optimal}. It applies only this correction
directly to raw CSI, omits SHARP's sparse path estimation, and then uses the
fixed Doppler transform.

This baseline matches the classifier accuracy of the official precomputed
SHARP Doppler at \(91.36\%\). Their individual predictions agree on \(98.33\%\) of
windows. Thus, the methods make different mistakes on a small number of samples, and these differences cancel out, resulting in the same overall accuracy. The affine
method obtains off-center MSE \(1.55\times10^{-4}\), F1 \(0.943\), and
motion-mass ratio 0.982.

This establishes that the local raw windows contain enough information for a
high-utility representation. It also highlights the main remaining neural error:
the direct model learns the approximate timing and location of motion despite unwanted phase noise. However, its generic decoder often replaces SHARP’s sharp, sample-specific peaks with smoother average patterns.
Under this setup, explicit affine phase modeling is highly effective and leaves
a substantially smaller gap than the corrected generic neural decoder.

\subsection{Limitations and Answer to the RQ}

The experiments use one seed, one frozen classifier, and 92 aligned AR
recordings paired with official precomputed SHARP Doppler traces; five candidate pairs are
excluded for temporal mismatch. The classifier test measures compatibility
with SHARP's established interface, not
the performance of a newly trained classifier adapted to student maps. 

The answer to the primary question is therefore mixed but operationally
negative. Direct regression can learn a meaningful approximation: the corrected
model generalizes event support and reaches \(63.09\%\), so the task is not
infeasible. It does not preserve enough fine structure to replace SHARP at
\(91.36\%\), especially for walking and jumping. The most promising next step is a phase-aware hybrid model that first cleans the CSI and then learns only the remaining correction or transformation.


\section{Conclusion}
\label{sec:conclusion}

Direct raw-CSI regression learned meaningful held-out motion support and the
complete corrected configuration improved over the previous neural baseline,
but it was not classifier-equivalent to the official precomputed SHARP Doppler
representation.
Its main failure was spectral and temporal smoothing, especially for walking
and jumping. Explicit affine phase correction followed by the fixed Doppler
transform performed substantially better under the tested setup. This supports
future phase-aware or hybrid learned models, such as learning after explicit
sanitization or unrolling the sparse phase stage. Because architecture,
normalization, alignment, sampling, batch size, and schedule changed together,
the experiments do not isolate each contribution; repeated seeds and equivalent
runtime measurements also remain future work.


\section{Exam rules}

What you need to do to pass the exam:
\begin{itemize}
\item Optional: team up with other students. Max. group size is \textbf{three students} per group;
\item Identify a project to work on, devise your own neural network architecture and test it on the provided dataset;
\item \textbf{Pass the written exam:} You can start working on the project before passing the written exam but you can only submit the report and code after passing it.
\item \textbf{Prepare a written project report:} Follow the template above to be sure to include all the required content;
\item \textbf{Submit the project for evaluation:} Submit i) your written report and ii) the code \textit{as two separate files};
\end{itemize}

Your final project score will be obtained taking into account the following criteria:
\begin{itemize} 
\item \textbf{Project} (70 points): originality (20 pt.) - data preprocessing techniques (10 pt.) - learning architectures (20 pt.) - comparison against baseline/other approaches (20 pt.) 
\item \textbf{Written report} (30 points): clarity (10 pt.) - completeness (10 pt.) - analysis of results (number and type of metrics used) (10 pt.)
\end{itemize}

The score on the project is computed as
\begin{equation}
\textrm{Project score} = \frac{\textrm{tot\_points} \times 30}{90}.
\end{equation}

The final grade will be computed as the average of the score in the written exam (50\%) and that on the project (50\%).


\bibliography{biblio}
\bibliographystyle{ieeetr}

\end{document}


